\chapter{How Much Compute Does a Nation's Science Need?}
\chaptermark{Sizing Sovereign AI4S}
\label{ch:sizing}

Every argument in this book has so far been about who supplies capability and at what price. This chapter asks the reciprocal question, which almost no government has answered in public: \emph{how much does a country actually need?} It matters here for three reasons. It converts the sovereignty argument from a posture into a number, and numbers can be checked. It exposes---quantitatively---the gap between what nations are announcing and what their own research populations would consume, which turns out to be an order of magnitude or more. And it supplies the demand side of the institutional proposal in Chapter~\ref{ch:consortium}: one cannot design a federation to pool capacity without a shared, comparable estimate of the capacity required. The material here follows a measurement-anchored sizing framework developed at RIKEN R-CCS and its three national instantiations~\cite{ai4s-framework}; the epistemic discipline is the framework's own, and it matches this book's: every input is typed by what would make it trustworthy.

\begin{quote}\small
\textbf{Why this chapter is in the main text.} A sizing framework detailed enough to be checked can read as a second book inserted into the first, and the temptation is to relegate it to an appendix. The temptation is resisted, and the reason is Chapter~\ref{ch:prize}: if the largest value in AI is scientific, then the question ``how much compute does a nation's science need, and what does it cost to rent instead of own?'' is not a technical annexe to the convergence thesis---it is the thesis, evaluated on the workload that matters most. What the appendix instinct is right about is the connective tissue, and this chapter supplies it: every section below closes by naming the chapter it feeds.
\end{quote}

\section{The measurement gap}

National AI-for-Science infrastructure has, until recently, been sized without measurement. Procurement scales were set by budget envelopes, vendor availability, flagship-application extrapolation, or elicitation workshops---none of which measures the per-researcher demand of the workload class the infrastructure is being built for, namely LLM agents orchestrating simulation, data analysis, and autonomous experimentation~\cite{ai4s-framework}. Meanwhile the deployment wave is global and simultaneous: the US Genesis Mission systems, EuroHPC's AI factories, the UK's Isambard-AI, Japan's Rikyu and FugakuNEXT line, and Singapore's national AI programs are all sizing sovereign AI4S capacity now, mostly without a common yardstick.

Three developments make a portable framework possible. The demand side acquired a measured anchor: a 56-day operational record of an early-access agentic campaign on a GB200-class system---101 vetted users, 53,173 jobs, 340,397 GPU-hours, saturating an $\sim$800-GPU partition within five weeks---which yields a revealed-demand floor of 443 delivered BGPU-hours per active heavy-user day that any country can adopt as a starting scenario and test against its own probe. The supply side converged onto public, reproducible serving benchmarks. And the policy context became genuinely multi-national, so that comparability now has customers~\cite{ai4s-framework}. What the OECD's Frascati Manual did for R\&D statistics---not identical answers across countries, but commensurable ones---is what a shared, openly gated estimation framework could do for research-compute demand.

\section{A vendor-neutral unit, and a conversion rule}

\begin{quote}\small
\textbf{Definition (Baseline GPU).} One BGPU is the AI-serving throughput of a reference accelerator possessing an 8\,TB/s HBM memory system---concretely a B200/B300-class device. It is a \emph{planning unit}, not a product: any accelerator converts into it by the bandwidth-first, benchmark-adjusted rule below, with a stated adjustment band. A separate and never-summed unit, BGPU$_{\mathrm{BW}}$, denotes 8\,TB/s of delivered bandwidth for classical modelling and simulation. Conversion uncertainty is carried explicitly as the $\delta$ band of Table~\ref{tab:bgpu} and propagates into every capacity figure in this chapter; a $\pm$10\% band on $\delta$ moves national requirements by the same proportion, which is smaller than the population-definition sensitivity discussed below and much smaller than the behavioural-transfer uncertainty.
\end{quote}

All capacities and demands are stated in \textbf{Baseline GPUs} (BGPU): the AI-serving throughput of one NVIDIA B200/B300-class accelerator---concretely, a device with an 8\,TB/s HBM memory system, the resource that governs long-context decode. The unit is deliberately not a product. Any accelerator $g$ converts by a bandwidth-first, benchmark-adjusted rule
\[
\gamma_g \;=\; \frac{\mathrm{BW}_g}{8\ \mathrm{TB/s}} \times \delta_g ,
\]
where the bandwidth term is physics and $\delta_g$ absorbs compute-side and software-stack differences, estimated from published matched-condition benchmark results and carrying a stated band~\cite{ai4s-framework}. The rule's worked example matters politically as well as technically: AMD's MI355X carries 288\,GB of HBM3E at 8.0\,TB/s---a bandwidth ratio of exactly 1.00---and its published MLPerf record supports $\delta\approx1$, so it converts at 1.0 BGPU, largely equivalent to Blackwell. \emph{A requirement denominated in BGPU is procurable from any vendor whose devices convert at a defensible factor}, which is precisely what a public planning framework must guarantee and what a vendor-named unit cannot. Table~\ref{tab:bgpu} gives the adopted factors.

\begin{table}[htbp]
  \centering
  \small
  \caption{Cross-vendor BGPU conversion: HBM bandwidth ratio, benchmark-adjustment band $\delta$, and adopted planning factor $\gamma$~\cite{ai4s-framework}. Scenario entries lack a public reference-class serving result.}
  \label{tab:bgpu}
  \begin{tabular}{llrrl}
    \toprule
    Device & Vendor & BW (TB/s) & $\gamma$ (BGPU) & Basis \\
    \midrule
    H100 & NVIDIA & 3.35 & 0.5 & $\delta$ 1.0--1.2 \\
    H200 & NVIDIA & 4.8 & 0.5 & $\delta$ 0.8--1.0 \\
    \textbf{B200 / GB200 / B300} & NVIDIA & 8.0 & \textbf{1.0} & definitional \\
    MI300X & AMD & 5.3 & 0.6 & $\delta$ 0.85--1.0 \\
    \textbf{MI355X} & AMD & 8.0 & \textbf{1.0} & $\delta$ 0.9--1.1 (MLPerf-supported) \\
    TPUv7 (Ironwood) & Google & 7.37 & $\sim$0.9 & scenario \\
    R200 (Rubin) & NVIDIA & 22.0 & 2.5 & $\delta$ 0.91 derate \\
    F200-class (2030 flagship) & --- & $\sim$48 & 6 (rounded) & \emph{derived}, not asserted \\
    \bottomrule
  \end{tabular}
\end{table}

The 2030 flagship-generation factor of 6 is derived rather than taken from a vendor roadmap: dividing a published program disclosure of $\geq$800\,PB/s aggregate accelerator memory bandwidth over a disclosed fleet of $\sim$15,000--20,000 devices gives 40--53\,TB/s per device, i.e.\ 5.0--6.7$\times$ the 8\,TB/s baseline; the planning value is rounded to 6, and because the bandwidth target is a floor, these are lower bounds~\cite{ai4s-framework}. Installed fleets deliver $H_{\mathrm{yr}} = \alpha\times 8{,}760 = 7{,}884$ effective GPU-hours per GPU-year at fleet availability $\alpha=0.90$. Classical modeling-and-simulation demand, where a country carries it, is expressed in a separate bandwidth currency (BGPU$_{\mathrm{BW}}$, one unit $=$ 8\,TB/s of delivered bandwidth) and is \emph{never} summed with serving-BGPU---a discipline that prevents the most common category error in national compute accounting.

\subsection{From scalar to workload-weighted vector}
\label{sec:bgpu-vector}

The framework owes one generalization an explicit statement, and this subsection supplies it. A bandwidth-first scalar is defensible because long-context, memory-bound decode dominates the measured national AI4S workload---but the workload is a portfolio, and a device strong on decode can be weak on training software, network injection, or reliability. The disciplined form defines a per-device capability vector,
\[
\mathbf{B}_g \;=\; \left( B_{\mathrm{decode}},\; B_{\mathrm{prefill}},\; B_{\mathrm{training}},\; B_{\mathrm{simulation}},\; B_{\mathrm{network}},\; B_{\mathrm{storage}},\; A_{\mathrm{availability}} \right),
\]
where each of the first six components is the device's sustained throughput on that workload class \emph{divided by} the reference device's sustained throughput on the same class under the matched-condition protocol of Section~\ref{sec:archtheory}---so each is a dimensionless ratio equal to 1.0 for the reference device---and $A_{\mathrm{availability}}$ is the fraction of wall-clock time the device is usable in production, likewise normalized to the reference fleet's $\alpha=0.90$. The scalar planning unit is derived only \emph{after} applying a declared national workload mix $\{\omega_k\}$, non-negative weights summing to one over the seven components:
\[
\mathrm{BGPU}_{\mathrm{plan}}(g) \;=\; \sum_k \omega_k\, B_{g,k}.
\]
The two-layer demand model already supplies the mix: the token-service layer weights decode and prefill; the scientific-execution layer weights training, simulation, network, and storage; and the published national instantiations state their $\omega_k$ implicitly through the layer split. What the vector form adds is protection against exactly the failure that matters: a device with spectacular memory bandwidth but immature training software, poor injection bandwidth, or weak fleet reliability can no longer receive a flattering conversion factor, because its $B_{\mathrm{training}}$, $B_{\mathrm{network}}$, and $A_{\mathrm{availability}}$ components are carried explicitly and the declared mix prices them. The scalar survives for politics---ministries budget in one number---but the vector is now the definition and the scalar is its declared projection, with the mix printed beside every national requirement. (The LX2-class integrated processor is the instructive case: strong $B_{\mathrm{decode}}$ and $B_{\mathrm{simulation}}$, unproven $B_{\mathrm{training}}$, unknown $A_{\mathrm{availability}}$ at fleet scale---the vector says precisely what the proof suite of Section~\ref{sec:proofsuite} must fill in before a conversion factor is defensible.)

\section{The two-layer demand model}

A nation's requirement for target year $Y$ is the sum of two concurrently provisioned layers, $G_{\mathrm{nat}}(Y) = G_{\mathrm{tok}}(Y) + G_{\mathrm{sci}}(Y)$: a token-service layer and a scientific-execution layer (Figure~\ref{fig:ai4s-model}). Both must be provisioned simultaneously, because the agentic loop does not close otherwise---the orchestrating model is idle without the batch jobs it launches, and the batch jobs are unlaunched without the orchestrator.

\begin{figure}[htbp]
\centering
\begin{tikzpicture}[
  every node/.style={font=\footnotesize},
  layer/.style={draw=inkMut, line width=0.7pt, rounded corners=3pt, align=left,
                inner sep=5pt, text width=104mm},
  cell/.style={draw=inkMut, line width=0.5pt, rounded corners=2pt, fill=white,
               align=center, font=\scriptsize, inner sep=3pt, text width=44mm},
  arr/.style={-{Stealth[length=5pt]}, line width=1pt, color=inkSec},
]
% --- token layer
\node[layer, fill=sOne!8, draw=sOne] (tokhdr) at (0,0)
  {\textbf{Token-service layer}\hfill $G_{\mathrm{tok}} = T_{\mathrm{nat}}\cdot k$\\[1pt]
   \scriptsize interactive serving, SLO-provisioned};
\node[cell, fill=sOne!5, below left=3mm and 0mm of tokhdr.south, anchor=north east, xshift=-2mm] (master)
  {\textbf{master}: 1T-class MLA MoE\\ $L_m = 2^{20}$ ctx, $\Theta_m = 115$ tok/s\\
   \textbf{10\% of tokens, 88.7\% of cost}};
\node[cell, fill=sOne!5, below right=3mm and 0mm of tokhdr.south, anchor=north west, xshift=2mm] (sub)
  {\textbf{sub-agents}: 120B SWA\\ 128K ctx, $\Theta_s = 8{,}161$ tok/s\\
   90\% of tokens, 11.3\% of cost};

% --- execution layer
\node[layer, fill=sTwo!8, draw=sTwo, below=36mm of tokhdr] (scihdr)
  {\textbf{Scientific-execution layer}\hfill
   $G_{\mathrm{sci}} = G^{\mathrm{AI}}_{\mathrm{sim}} + G_{\mathrm{tra}} + G_{\mathrm{dat}}$\\[1pt]
   \scriptsize batch partition and resident pipelines};
\node[cell, fill=sTwo!5, below left=3mm and 0mm of scihdr.south, anchor=north east, xshift=-2mm] (sim)
  {simulation: MD, DFT, CFD,\\ plus persistent surrogate programs};
\node[cell, fill=sTwo!5, below right=3mm and 0mm of scihdr.south, anchor=north west, xshift=2mm] (tra)
  {training \& fine-tuning\\ (GPU-substrate-pinned)};
\node[cell, fill=sTwo!5, below=3mm of scihdr.south, anchor=north, yshift=-14mm, text width=94mm] (dat)
  {facility pipelines: resident streaming inference at detector rate,
   $G_{\mathrm{dat}} = F_{\mathrm{eq}}\,g_{\mathrm{fac}}$};

% --- coupling arrows, drawn in the clear vertical gap
\draw[arr] ($(master.south)+(-14mm,-3mm)$) -- ++(0,-13mm)
  node[midway, left, font=\scriptsize\color{inkSec}, align=right, xshift=-1mm]
  {agent-launched\\batch jobs};
\draw[arr] ($(sub.south)+(14mm,-16mm)$) -- ++(0,13mm)
  node[midway, right, font=\scriptsize\color{inkSec}, align=left, xshift=1mm]
  {results return\\to context};

\node[font=\scriptsize\color{inkSec}, align=center, below=4mm of dat, text width=104mm]
  {$G_{\mathrm{nat}}(Y)=G_{\mathrm{tok}}(Y)+G_{\mathrm{sci}}(Y)$ --- both layers must be provisioned
   \emph{concurrently} for the agentic loop to close};
\end{tikzpicture}
\caption[The two-layer national AI4S demand model]{The demand model. A nation's agentic AI-for-Science
requirement is the concurrent sum of a token-service layer, priced through a canonical master/sub-agent serving
configuration, and a scientific-execution layer decomposed by function into simulation, training, and data
processing---with training pinned to GPU-class substrate and simulation the only cross-substrate routable
term. Note the master asymmetry carried in the top row: a tenth of the tokens, generated at million-token
context, consume nearly nine-tenths of the token-service capacity. After the sizing
framework~\cite{ai4s-framework}; the serving configuration is a versioned reference object [U]-r/[P], not a
physical constant.}
\label{fig:ai4s-model}
\end{figure}


\textbf{The token layer.} The open-tier population $N$ (the number of researchers entitled to draw on the national service) is partitioned into service tiers indexed by $i$, with population shares $\varphi_i$ (summing to one), duty factors $u_i$ (the fraction of the year a researcher in tier $i$ is actively drawing tokens rather than idle), and sustained output-token rates $r_i$ while active; the portable default is a three-tier profile of 10\%/20\%/70\% of researchers at 2,000/400/150 Mtok per user-month. The national output stream is $T_{\mathrm{nat}} = N\sum_i \varphi_i\, u_i\, r_i$ tokens per second, and its serving cost is $G_{\mathrm{tok}} = T_{\mathrm{nat}}\cdot k$, where $k$ is the BGPU cost per token per second of a canonical, versioned agentic serving configuration: a \emph{master} model---a trillion-parameter MLA mixture-of-experts holding a $2^{20}$-token campaign context---delegating to cheap, high-throughput 120B-class sliding-window \emph{sub-agents} at 128K context, with 10\% of national tokens going to the master~\cite{ai4s-framework}.

That configuration produces the single most consequential structural fact in the framework, and Figure~\ref{fig:ai4s-plots} (top) states it: at $k = 9.79\times10^{-4}$ BGPU per token per second, \textbf{88.7\% of the token-service capacity is the master term}. Ten percent of the tokens, generated at million-token context, cost nine-tenths of the token-service GPUs. The reason is mechanical. Decode throughput on one accelerator is delivered memory bandwidth divided by the bytes that must be read to emit one token, and at long context that traffic is dominated by the key--value cache. Writing $\mathrm{BW}$ for delivered bandwidth (8\,TB/s at the BGPU reference), $L_m$ for the master's resident context length in tokens, and $b_{\mathrm{KV}}$ for the KV-cache bytes per token of context, master throughput is
\[
\Theta_m \;=\; \frac{\mathrm{BW}}{L_m\, b_{\mathrm{KV}}} \;=\; 115 \ \text{tokens per second},
\]
against $\Theta_s = 8{,}161$ for the 128K-context sub-agents---a 71-fold ratio driven entirely by context length against KV traffic. That ratio has a direct architectural consequence for this book: \emph{national AI4S demand is overwhelmingly a long-context memory-capacity and memory-bandwidth problem}, which is exactly the workload the capacity-first architecture of Chapter~\ref{ch:compute} is built for and exactly the workload on which token pricing is most punishing.

\begin{figure}[htbp]
\centering
\begin{tikzpicture}
\begin{axis}[bookbar, height=4.3cm, width=0.66\linewidth, bar width=15pt,
  ylabel={Share (\%)},
  title={The master's asymmetry},
  symbolic x coords={{tokens generated},{token-service GPUs}},
  xtick=data, ymin=0, ymax=115,
  xticklabel style={font=\scriptsize\color{inkSec}},
  nodes near coords, nodes near coords style={font=\scriptsize\color{inkSec}},
  legend style={at={(0.5,-0.30)}, anchor=north, legend columns=2},
]
\addplot[fill=sOne] coordinates {({tokens generated},10) ({token-service GPUs},88.7)};
\addlegendentry{master (1M-token context)}
\addplot[fill=sThree] coordinates {({tokens generated},90) ({token-service GPUs},11.3)};
\addlegendentry{sub-agents (128K context)}
\end{axis}
\end{tikzpicture}

\vspace{3mm}

\begin{tikzpicture}
\begin{axis}[bookaxis, height=5.4cm,
  ymode=log,
  ylabel={Capacity (BGPU, log scale)},
  title={The universal-reach ladder: committed against required},
  symbolic x coords={Singapore,Japan,{United States}},
  xtick=data, ymin=300, ymax=3e6,
  ytick={1000,10000,100000,1000000},
  yticklabels={$10^3$,$10^4$,$10^5$,$10^6$},
  xticklabel style={font=\small\color{inkSec}},
  ybar, bar width=7pt, xmajorgrids=false, enlarge x limits=0.36,
  every axis plot/.append style={draw=none},
  legend style={at={(0.5,-0.22)}, anchor=north, legend columns=2},
]
\addplot[fill=sOne] coordinates {(Singapore,1000) (Japan,39000) ({United States},112000)};
\addlegendentry{committed fleet}
\addplot[fill=sThree] coordinates {(Singapore,3200) (Japan,28300) ({United States},62300)};
\addlegendentry{rung 1: uniform baseline service}
\addplot[fill=sTwo] coordinates {(Singapore,8300) (Japan,72700) ({United States},160000)};
\addlegendentry{rung 2: full-population tiered service}
\addplot[fill=sFour] coordinates {(Singapore,57000) (Japan,503000) ({United States},1130000)};
\addlegendentry{modeled national requirement}
\end{axis}
\end{tikzpicture}

\vspace{2mm}

\begin{tikzpicture}
\begin{axis}[bookbar, height=4.2cm, width=0.72\linewidth, bar width=15pt,
  ylabel={BGPU per open-tier researcher},
  title={Provisioning against requirement, per researcher},
  symbolic x coords={Singapore,Japan,{United States},{requirement}},
  xtick=data, ymin=0, ymax=1.15,
  xticklabel style={font=\scriptsize\color{inkSec}},
  nodes near coords, nodes near coords style={font=\scriptsize\color{inkSec}},
  every node near coord/.append style={/pgf/number format/fixed,
     /pgf/number format/precision=2, /pgf/number format/fixed zerofill},
]
\addplot[fill=sOne] coordinates
  {(Singapore,0.02) (Japan,0.08) ({United States},0.10) ({requirement},1.00)};
\end{axis}
\end{tikzpicture}
\caption[The AI4S sizing results]{Three readings of the same framework~\cite{ai4s-framework}. \emph{Top}: ten
percent of tokens, generated at million-token context, consume nearly ninety percent of the token-service
capacity---the configuration's central structural fact, invariant across countries, and the reason master context
length is the layer's primary scenario axis. \emph{Middle}: committed national fleets against two priced service
rungs and the full modeled requirement, log scale. Japan's and the United States' commitments already exceed the
uniform-baseline rung but reach only $\sim$54\% and $\sim$70\% of the tiered rung; Singapore's clears neither.
\emph{Bottom}: the same result per researcher---announced infrastructures provision 0.02--0.10 BGPU against a
modeled requirement near 1.0, an underprovision of roughly one order of magnitude for Japan and the US and
closer to fifty-fold for Singapore. Japan is measurement-anchored [M/V]; Singapore and the US are
pre-measurement scenarios [T] whose behavioral parameters are transferred from the Japanese probe and are
labeled as such.}
\label{fig:ai4s-plots}
\end{figure}


\textbf{The execution layer.} Defined function-first and estimated operator-first:
\[
G_{\mathrm{sci}} = \underbrace{G^{\mathrm{AI}}_{\mathrm{sim}}}_{\text{simulation}} + \underbrace{G_{\mathrm{tra}}}_{\text{training}} + \underbrace{G_{\mathrm{dat}}}_{\text{data processing}},
\qquad
W_{\mathrm{ag}} = N\cdot f_{\mathrm{heavy}}\cdot d_{\mathrm{duty}}\cdot w_{\mathrm{heavy}}\cdot 365 ,
\]
with $G^{\mathrm{AI}}_{\mathrm{sim}} = (\chi_{\mathrm{sim}}W_{\mathrm{ag}} + n_{\mathrm{prog}}w_{\mathrm{prog}})/H_{\mathrm{yr}}$, $G_{\mathrm{tra}} = (1-\chi_{\mathrm{sim}})W_{\mathrm{ag}}/H_{\mathrm{yr}}$, and $G_{\mathrm{dat}} = F_{\mathrm{eq}}\,g_{\mathrm{fac}}$.

Every symbol, in order. $W_{\mathrm{ag}}$ is the annual agentic-execution workload in \emph{delivered BGPU-hours}: the heavy-user fraction $f_{\mathrm{heavy}}$ (default 0.125) times the activity-by-discipline duty product $d_{\mathrm{duty}}$ (0.31---the share of days a heavy user is actually active, folded with the share of disciplines that draw on this workload) times $w_{\mathrm{heavy}}$, the measured per-heavy-user daily driver (443 delivered BGPU-hours per active claimant-day), times 365 days, times the population $N$. Dividing BGPU-hours by $H_{\mathrm{yr}}$---the 7,884 effective hours one installed BGPU delivers per year---converts a flow of work into a stock of machines, which is why $H_{\mathrm{yr}}$ appears in the denominator of the first two terms. The superscript AI on $G^{\mathrm{AI}}_{\mathrm{sim}}$ marks AI-driven simulation specifically, distinguishing it from the classical modelling-and-simulation demand that is accounted separately in BGPU$_{\mathrm{BW}}$ and never summed with it. $\chi_{\mathrm{sim}} \in [0,1]$ is the \emph{simulation share}: the fraction of agentic execution that is surrogate-and-simulation work rather than model training, so $\chi_{\mathrm{sim}}$ and $(1-\chi_{\mathrm{sim}})$ split $W_{\mathrm{ag}}$ between the first two terms. $n_{\mathrm{prog}}$ is the census of persistent-surrogate programs and $w_{\mathrm{prog}}$ their annual cost, 0.1--0.5M BGPU-hours each. $F_{\mathrm{eq}}$ is the national inventory of large-instrument facility equivalents and $g_{\mathrm{fac}}$ the installed BGPU each requires for its data-processing pipeline, 2,400--3,250 apiece~\cite{ai4s-framework}. An explicit substrate-eligibility matrix pins training to GPU-class systems and makes simulation the only cross-substrate routable function---which is where a country's classical fleet can legitimately offset AI demand, and where it cannot.

\section{Typed parameters and transfer rules}

The framework's portable core is not the arithmetic but the \emph{parameter taxonomy}, and it is the closest methodological cousin this book has found to its own evidence grading. (One notational warning before the classes are named: this taxonomy's labels are \emph{local to this chapter} and are not the book's evidence grades of Section~\ref{sec:evidence}. In particular \textbf{[M]} here means ``must-measure,'' not ``author-modeled,'' and \textbf{[P]} here means ``policy choice,'' not ``primary-source claim.'' The two systems answer different questions---the global grades ask how well a claim is evidenced, the parameter classes ask whether a number may be carried across a border---and where both appear on a page, the parameter classes are the ones attached to model \emph{inputs}.) Every input is classed \textbf{[U]} universal (serving physics and benchmark calibrations that transfer across borders), \textbf{[M]} must-measure (populations, facility inventories, electricity tariffs---country-specific by construction), \textbf{[T]} transferred (behavioral quantities measured so far only in one country, carried elsewhere as labeled scenario values until replaced by a domestic probe), or \textbf{[P]} policy (allocation choices, which are decisions rather than measurements), with an orthogonal \textbf{[F]} future-scenario modifier for forward technology assumptions~\cite{ai4s-framework}. The taxonomy states which numbers may cross a border, which must never, and which may travel only with a warning label and an expiry date.

That last category is where intellectual honesty lives. The serving physics---memory-bandwidth-limited decode of trillion-parameter MoE models, calibrated delivered efficiency, generation conversion factors---transfers exactly, because benchmarks do. What a heavy national research cohort actually does under scarcity transfers only as a hypothesis: Singaporean or American researchers may behave like the measured Japanese cohort, but no one has yet checked. A portable framework must therefore do more than parameterize the arithmetic; it must \emph{label every input by what would make it trustworthy, and specify the experiment that replaces each assumption with a local measurement}. The seven-step measurement protocol does exactly that---population construction, facility inventory, a saturating demand probe on a $\sim$1,000-GPU current-generation partition with full accounting, serving-basis calibration, classical-mission accounting, cost localization, and gated publication---so that ``run this analysis for country~X'' becomes an executable checklist rather than a research project.

\section{What three countries find}

The framework has been instantiated at both extremes of national scale (Figure~\ref{fig:ai4s-plots}). \emph{Japan}, the fully worked measurement-anchored example: $N\approx 500{,}000$ open-tier researchers, central 2030 requirement 497--510k BGPU. \emph{Singapore}, a pre-measurement scenario: $N\approx 57$k researchers, 55--59k BGPU---the entire national requirement smaller than Japan's proposed first-stage procurement, with the existing $\sim$1,900-GPU national research fleet at 1.5--2\% workload coverage and a proposed 3,000-GPU Rubin-class first stage reaching 14--16\% coverage and $\sim$64\% population reach for a four-year subtotal of \$0.39--0.44B. \emph{United States}: $N\approx 1.1$M, 1.09--1.17M BGPU central, with the committed Equinox-plus-Solstice systems---the largest AI procurement in DOE history---standing at $\sim$10\% workload coverage of the modeled requirement, closely paralleling Japan's no-new-procurement endpoint~\cite{ai4s-framework}.

Four disciplines govern how those numbers should be read, and the third and fourth are corrections to the loose way such figures are ordinarily handled. \emph{Installed peak, deliverable annual service, SLO-provisioned capacity, and effective utilization are four different quantities}, and a comparison that mixes them is meaningless; this chapter quotes the analytical design point and the operationally provisioned point separately wherever both exist. \emph{Estimands must match}: delivered usage $\ne$ requested work $\ne$ latent demand $\ne$ installed capacity $\ne$ population reach $\ne$ workload coverage. \emph{The scale-invariance finding is a hypothesis, not a result}---it holds across three instantiations of which only one is measured, so the honest statement is that the per-researcher core is \emph{modelled} as scale-invariant because both of its terms scale linearly in $N$ at fixed behaviour, and testing that assumption outside Japan is the first job of any second national probe. And \emph{the Singapore and US figures should not be quoted with more precision than transferred behavioural assumptions permit}: a reasonable confidence band on the $\approx$0.94 BGPU-per-researcher core is at least $\pm$40\% before national measurement, which is why this chapter's strategic conclusion is stated as an order-of-magnitude gap rather than as a percentage.

One further sensitivity deserves its own line, because it cuts against the chapter's own conclusion. \emph{Cheaper models may not reduce BGPU demand at all.} If capability per token rises and price per token falls, the rational institutional response is to run more agents, longer contexts, and deeper search---the Jevons argument of Section~\ref{sec:jevons} applied to science. Whether efficiency reduces or induces national compute demand is an empirical question that the measurement protocol can answer and has not yet; both signs are consistent with the framework, and this book's own bet, given the master-term arithmetic and the agentic direction of scientific workflows, is on induction rather than reduction.

Three structural findings survive across the $\sim$20$\times$ range of national scale. \emph{The per-researcher core is modelled as scale-invariant}: the token layer plus the agentic execution term compose to $\approx$0.94 BGPU per open-tier researcher in every country, because both scale linearly in $N$ at fixed behavior. \emph{The deviations are diagnostic}: what does not scale with $N$ is the facility-and-program term, set by instrument bandwidth and pipeline depth---$\sim$5.5\% of the total for Singapore, $\sim$9\% for the United States, where the streaming-inference class alone runs 60--132k BGPU, between half and 1.2 times the entire committed AI fleet. Sizing exercises that count only central systems will structurally undercount this class. \emph{And the headcount rule is worth a factor of six}: an OECD-consistent FTE count of US open-sector researchers is $\approx$185k against an eligibility headcount of $\approx$1.1M---a definitional choice that swamps every other sensitivity, and therefore one that any official exercise must state, defend, and apply consistently.

\section{The universal-reach ladder, and the order-of-magnitude gap}

The framework's sharpest instrument is a computed \emph{universal-reach ladder}: the installed capacity at which 100\% of researchers are served, at each of two precisely labeled service rungs. Rung one, \emph{uniform baseline service}---every researcher receiving the baseline 150\,Mtok/month tier---costs 3.2k (Singapore) / 28.3k (Japan) / 62.3k (US) BGPU installed. Rung two, \emph{full-population tiered-token service} at the differentiated 10/20/70 entitlement profile, costs 8.3k / 72.7k / 160.0k~\cite{ai4s-framework}.

The distinction sharpens both readings, and Figure~\ref{fig:ai4s-plots} (middle) shows why. Japan's and the United States' committed fleets already \emph{exceed} the uniform-baseline rung---rung coverage 137\% and 180\% at the annual-average convention, 106\% and 139\% even at SLO provisioning---yet stand at only $\sim$54\% and $\sim$70\% of the tiered rung, while Singapore's existing fleet covers under a third of even the baseline. Closing the US tiered-service gap takes roughly 19,000 Rubin-class GPUs beyond current commitments. And the full requirement---tiered tokens plus full heavy compute---costs a further $\sim$6$\times$ again: 54k / 470k / 1,034k BGPU.

Stated per researcher (Figure~\ref{fig:ai4s-plots}, bottom), the result is the chapter's headline: \textbf{announced national infrastructures provision 0.02--0.10 BGPU per open-tier researcher against a modeled requirement near 1.0}---an underprovision of roughly tenfold for the United States and for Japan before its proposed first stage, and closer to fiftyfold for Singapore. One more asymmetry is worth recording because it is the reason this chapter exists: on the public record, Japan's stage-1 proposal is sized against a measured national demand model but is unfunded, while the US commitment is funded, contracted, and the largest in the world but is not yet sized against its research population. \emph{One side has the measurement without the money; the other has the money without the measurement.}

\section{What this means for the argument of this book}

Four consequences, each connecting to a thread already running.

\emph{The token price is the science budget.} Chapter~\ref{ch:consortium} derives, from a workload share of $\sim$14\% and a 20--60$\times$ purchased-token premium, that frontier tokens consume 76--91\% of an AI4S budget (Figure~\ref{fig:ai4s-cost}). This chapter supplies the mechanism underneath that estimate: the master term is 88.7\% of token-service capacity, so the expensive layer is precisely the long-context orchestration that no research workflow can avoid and that commercial pricing charges most for. A nation that rents its master layer has rented the majority of its research computing.

\emph{The requirement is a memory problem, and therefore an architecture problem.} If national demand is dominated by million-token-context decode over trillion-parameter sparse models, then the machine that serves it best is capacity-rich and bandwidth-balanced---the design point of Section~\ref{sec:lpddr6} and the LX2 class of Section~\ref{sec:lineshine}, not a dense-FLOP tensor part. The sizing framework and the architecture argument are the same argument in different units.

\emph{Sovereignty has a computable price, and it is high but not absurd.} At delivered B200-hour costs of \$1.84--2.19 across the three countries~\cite{ai4s-framework}, the gap between announced and required capacity is a national-budget question rather than a technological impossibility---which is exactly the condition under which pooling becomes rational. A federation that shares capacity across time zones and campaign cycles converts three separately inadequate national fleets into one adequate pooled one; that is the operational case for Chapter~\ref{ch:consortium}, and it now has a denominator.

\emph{And comparability is itself strategic.} A country that cannot state its requirement in a unit others recognize cannot join a federation, cannot audit a vendor's claim, and cannot tell whether a foreign price list is cheap or ruinous. The framework's real product is not the Japanese number; it is the shared definition---which is why the invitation it extends, to run the protocol and publish the parameter file, is the least expensive strategic action available to any government reading this book.
